\chapter{Introduction}

\section{Project Overview and Motivation}

This report documents the design and construction of the \textbf{AI-Driven Smart
Home Energy Optimizer} (hereafter \textbf{SHEO}). SHEO models a single
apartment building and helps each unit's resident understand and reduce its electricity
bill, while the building owner oversees the building as a whole. It performs four closely
related tasks. First, it \emph{monitors} the real-time power and cumulative energy of every
device in a unit. Second, it offers \emph{capability-driven control}: rather than hard-coding a
separate screen for each appliance, the system describes each device type by a
declarative schema of its controls and telemetry, and the same interface renders and
validates commands for a plug, a bulb, a fan, an air-conditioner, or a sensor without
per-device code. Third, it \emph{learns daily habits} from the observed usage and
proposes human-readable automation rules of the form \textbf{WHEN \dots\ THEN \dots},
each of which the user can read, edit, accept, or dismiss. Fourth---and this is the
headline contribution that the client singled out---it \emph{estimates and displays the
Vietnamese-dong saving} that each rule is expected to deliver, shown \emph{before} the
rule is accepted, and then reports the saving actually accrued over the current billing
cycle.

The motivation is concrete and topical. Residential electricity costs in Vietnam follow
a tiered tariff in which the marginal price rises with consumption, so even small,
habitual waste compounds into a noticeably larger monthly bill. Residents are willing
to automate their appliances, but they are reluctant to trust automation they cannot
understand or whose benefit they cannot quantify. SHEO therefore treats
\emph{explainability and monetary visibility} as first-class requirements: every
suggestion is a sentence a non-expert can read, and every saving figure reduces to one
transparent formula expressed in VND. The project's purpose is to demonstrate that a
modest, fully working system, engineered with discipline, can deliver this value
end-to-end.

\section{Engineering Approach}

SHEO is, before anything else, an exercise in disciplined software engineering, and the
report is structured to make that discipline visible. Two principles govern it
throughout. First, the report follows the software lifecycle end to end---process and
management, requirements, design and construction, configuration management, and quality
assurance---so that every artefact has an explicit home. Second, it keeps a strict
separation between \emph{requirements} (\emph{what} the system must do) and \emph{design}
(\emph{how} it does it): the requirements chapter states only observable behaviour, while
the design and construction chapters describe the architecture, components, and code that
realise it. That separation is what keeps the chain from requirement to design to code to
test auditable.

\section{The Client and the Stakeholders}

Requirements work depends on a distinction that is easy to blur in everyday speech: the
\textbf{Client} (the party for whom the software is built and who funds it), the
\textbf{Customer} (the party who purchases or selects the software for use within an
organisation), and the \textbf{User} (the person who actually operates the software day
to day). Identifying these roles correctly is not a formality---it determines whose
concerns the design must privilege.

The roles for SHEO were fixed in the initial client session. The \textbf{Client is an
IoT contractor}: it manufactures and supplies the smart devices, builds the underlying
artificial-intelligence capability, and commissions this development team to build the
application that ties the two together. The Client funds the work and sets the
priorities but does not itself operate the running system; it is therefore a
\emph{business stakeholder outside the system boundary}, not a login actor. The
\textbf{Customer} is the \emph{building owner or landlord} that adopts SHEO
and runs it for a single apartment building; in the running system the Customer is realised as the
\textbf{Administrator} actor, who owns no unit of their own but oversees every unit---viewing
the resident roster, each unit's live power and usage, and the building totals---sells and
onboards units to residents (and may offboard a resident, returning the unit to vacant while
keeping its devices and history), and sets the building-wide tariff and billing cycle. The
Administrator never operates a unit's devices and authors no automation rules. The
\textbf{End-users} are the \textbf{Residents} (tenants), each occupying one unit, who view
their own unit's consumption and savings, control the operable devices in that unit, and
author their own automation rules, with no visibility of any other unit. Two further actors
complete the picture: a \textbf{Developer/Tester}, who prepares and reproduces the
demonstration, and a non-human \textbf{System Scheduler}---a time actor that
autonomously evaluates and fires rules and mines usage habits. The relationship between
the external stakeholders and the system is shown in Figure~\ref{fig:context}.

\fig{fig01-context}{System context overview. The three human actors---Administrator
(the Customer / building owner), Resident (tenant of a unit), and Developer/Tester---interact
with SHEO. The Client (IoT contractor) is a business stakeholder \emph{outside} the
boundary that commissions and funds the build. A mock simulator stands in for physical
devices, and the building-wide tariff is configured by the Administrator. The outputs returned to
users are the live dashboard, notifications, explainable rules, and VND savings.}{fig:context}

Table~\ref{tab:stakeholders} summarises each stakeholder, the role it plays with
respect to the system, and its principal concern. The most consequential entry is the
Customer's: in the client session the contractor assigned \emph{high priority} to the
requirement that the Customer be able to \emph{see the actual VND bill saved} by a
suggested rule. That single priority shapes much of the rest of the report---it drives
the bill-saving estimation engine, the explainability constraint, and the
headline-figure layout of the dashboard.

\begin{table}[htbp]
\centering
\caption{SHEO stakeholders, their roles with respect to the system, and their primary
concerns.}
\label{tab:stakeholders}
\begin{tabularx}{\linewidth}{L{2.6cm} L{4.4cm} Y}
\toprule
\textbf{Stakeholder} & \textbf{Role} & \textbf{Primary concern} \\
\midrule
Client (IoT contractor) & Commissions and funds the build; supplies devices and the
underlying AI. Business stakeholder, \emph{not} a system actor. & A demonstrable
product that proves its value, especially the monetary saving the Customer can see. \\
\addlinespace
Customer (building owner / landlord) & Adopts and operates SHEO; realised
as the \textbf{Administrator} actor, who owns no unit and oversees all of them. & A
building overview---the resident roster, every unit's usage, and building totals;
selling and onboarding units; and setting the building-wide tariff. \\
\addlinespace
End-user (Resident / tenant) & Day-to-day operator of one unit; the \textbf{Resident}
actor. & A clear at-a-glance view of the unit's own consumption and savings, and simple
control of that unit's operable devices. \\
\addlinespace
Developer / Tester & Prepares and reproduces the demonstration. & A deterministic,
reproducible system that is easy to exercise and test. \\
\addlinespace
System Scheduler & Non-human time actor inside the system. & Reliable, periodic
evaluation of rules and mining of habits. \\
\bottomrule
\end{tabularx}
\end{table}

\section{Scope and Objectives}

The objective of the project is to deliver, with full software-engineering rigour, the
five capabilities the client requested. These define the \emph{in-scope} functionality:
(1) real-time monitoring of device power and energy; (2) capability-driven device
control; (3) explainable WHEN--THEN automation rules with an opt-in auto-action;
(4) habit-learning recommendations derived from observed usage; and (5) estimation and
reporting of the VND bill saving of each rule. Each capability maps to a feature group
in the requirements specification and is traced through design, construction, and test.

Two deliberate exclusions place the work \emph{out of scope}. First, no physical IoT
hardware is used: real plugs, bulbs, fans, air-conditioners, and sensors are unavailable
to the team, so a \textbf{mock device simulator} stands in for them. The simulator
implements the same device interface that a real vendor integration would, so the
architecture remains hardware-ready, but the running product is demonstrable and
testable entirely in software. Second, the ``AI''---the habit detection the client (an
AI contractor) supplies---is isolated behind a swappable \textbf{recommendation-provider
port} (Strategy). The shipped default is an \textbf{explainable, rule-based} provider:
every detected habit surfaces as a rule the user can read. A black-box machine-learning
provider could be substituted behind the same port without touching the API, the UI, the
rule engine, or---decisively---the VND saving estimator, which is deliberately kept as
SHEO's own software so the saving always reduces to a documented formula. Shipping the
explainable provider by default keeps the product demonstrable in software and honours the
client's high-priority demand that the saving be \emph{visible and trustworthy}, while the
design admits a learned model where one is warranted.

Both decisions follow from a straightforward \emph{function / cost / time} trade-off.
Procuring and integrating real hardware would inflate cost and schedule for no gain in
demonstrated function, whereas a deterministic simulator delivers the same observable
behaviour. Likewise, training a model now would raise development cost and risk for no
gain in demonstrated function, whereas the default explainable provider delivers it
immediately---and the provider port keeps a learned model available without
re-architecting. In both cases the chosen scope maximises the function that matters while
respecting the cost and time available.

\section{Report Structure}

The remaining chapters follow the software-engineering lifecycle, tracing SHEO as a
single coherent system from problem to delivered, verified product.
Table~\ref{tab:roadmap} maps each chapter to its topic.

\begin{table}[htbp]
\centering
\caption{Report roadmap.}
\label{tab:roadmap}
\begin{tabularx}{\linewidth}{L{2.2cm} Y}
\toprule
\textbf{Chapter} & \textbf{Topic} \\
\midrule
Chapter~2 & Software development process and Agile practice. \\
Chapter~3 & Software project management. \\
Chapter~4 & Requirements analysis and modelling. \\
Chapter~5 & Software architecture and detailed design. \\
Chapter~6 & User-interface and user-experience design. \\
Chapter~7 & Software construction and code quality. \\
Chapter~8 & Software configuration management. \\
Chapter~9 & Quality assurance: verification, validation, and maintenance. \\
Chapter~10 & Conclusion and reflection. \\
\bottomrule
\end{tabularx}
\end{table}

Taken together, the chapters trace SHEO from the requirements that state \emph{what} it
must do, through the design and construction that determine \emph{how} it does it, to the
configuration management and quality assurance that keep it correct as it changes.
